\section{Index units and information boundaries}\label{sec:C16:index}
The indexing unit determines which relations retrieval can preserve. Fixed token chunks are simple but can separate an assertion from its qualification, a tool call from its result, or a correction from the statement it revises. Structural units such as sections, messages, episodes, and typed records can preserve those boundaries, but they vary in length and can contain multiple propositions with different authority. Chunking is part of the representation, not a neutral preprocessing detail.

A document index should retain source identity, version, location, timestamp, and access scope. An agent-memory index additionally needs event and session relationships. If a retrieved passage is regenerated from a summary, its identifier should distinguish the summary from the original source. Otherwise a citation can appear precise while pointing to an unsupported abstraction. A hierarchy can provide compact navigational text while keeping a path to exact evidence.

Dense retrieval uses learned representations to compare queries and candidates; lexical retrieval provides a complementary signal for exact strings, identifiers, and uncommon terms. Hybrid retrieval is an engineering response to different failure modes, not a universal guarantee. Dense similarity can miss a newly introduced identifier, while lexical overlap can retrieve a misleading near-duplicate. Query construction and metadata filtering should reflect the task rather than assuming one similarity function captures every relation.

Versioned indexing is necessary when documents change. A vector for an old passage can remain searchable after the displayed text has been replaced, creating an evidence mismatch. Store an index-generation identifier and source-content digest, and reject or refresh stale results according to the application contract. Re-embedding the entire corpus can be expensive; incremental updates require careful deletion and alias handling. The cost of these operations belongs in the lifecycle boundary.

Access filtering should occur before exposing content to an unauthorized model or worker. Filtering only the final answer does not undo the disclosure. Search ranking can itself reveal sensitive metadata if unauthorized records affect returned scores or counts. A strong application contract identifies which metadata may be visible, which store partitions are isolated, and how cached results are invalidated after access changes. Retrieval is an information-flow boundary as well as a ranking stage.

The September LGM method changes the representation boundary itself. It pools selected backbone hidden states, obtains sparse concept activations through an autoencoder, constructs query-dependent graph relations, and conditions generation through learned latent prefix tokens. Its training combines response generation, evidence reconstruction, and sparse reconstruction objectives~\cite{cumb2026lgmv2}. These are coupled learned mechanisms, unlike an external text graph used only for lookup. The ``neuro-symbolic'' designation does not establish logical soundness or uniquely identifiable concepts: sparse activations and learned message passing do not by themselves prove either. A structure ablation must therefore control backbone adaptation, prefix capacity, and evidence-reconstruction supervision before assigning the full gain to graph topology. Reusing the latent bank after changing its encoder or consumer also requires the compatibility tests of Chapter~\ref{ch:C17}.

A proposed index audit samples records across document formats and update histories, checks that every returned item resolves to the recorded version, and verifies that qualification and role boundaries survive chunking. It should include deleted sources, renamed artifacts, repeated identifiers, and contradictory revisions. These checks do not prove semantic usefulness, but they establish that the retriever is searching the evidence it claims to search.

